\documentclass[10pt,a4paper]{amsart}
\usepackage[margin=19mm]{geometry}
\usepackage{amsmath,amssymb,mathtools}
\usepackage[scr=rsfs,cal=euler]{mathalfa}
\usepackage{xfrac}
\usepackage[unicode,hidelinks,bookmarks=false]{hyperref}
\newcommand{\ie}{i.e.\ }
\newcommand{\cf}{cf.\ }
\newcommand{\resp}{respectively\ }
\newcommand{\Subsection}{\subsection}
\providecommand{\nobreakdash}{\nobreak\hbox{-}}
\setlength{\emergencystretch}{2em}
\multlinegap0pt
\usepackage{bm}
\usepackage{bbm}
\usepackage{dsfont}
\usepackage{xspace}
\usepackage{cancel}
\usepackage{mathtools}
\usepackage{amsthm}
\makeatletter
\@ifundefined{theorem}{\newtheorem{theorem}{Theorem}}{}
\@ifundefined{proposition}{\newtheorem{proposition}[theorem]{Proposition}}{}
\makeatother
\newcommand{\bbR}{\mathbb{R}}
\title[First-Order Analysis of Optimization in Uniformly Convex Met]{First-Order Analysis of Optimization in Uniformly Convex Metric Spaces: Directional Subderivatives and Basic Descent}
\author{D. Russell Luke, Titus Pinta, Qinyu Yan}
\date{Source: arXiv:2607.20677v1 (CC BY 4.0)}
\begin{document}
\maketitle

\noindent\emph{Source and licence.} First-Order Analysis of Optimization in Uniformly Convex Metric Spaces: Directional Subderivatives and Basic Descent. Version arXiv:2607.20677v1 (CC BY 4.0), distributed under the Creative Commons Attribution 4.0 International licence, as recorded in the arXiv licence metadata for this exact version and shown on the abstract page https://arxiv.org/abs/2607.20677v1. Full rights notice: licenses/arxiv-2607.20677-CC-BY-4.0.txt.\par\smallskip

\noindent\textit{Editorial scope.} Counted unit: the Proposition whose source label is \texttt{None} in the article (source0.tex lines 1235--1239, source label \texttt{None}), delivered together with the article's own complete original proof of it (source0.tex lines 1240--1332, one \texttt{proof} environment of 93 lines). No mathematics is written, repaired or re-derived: statement and proof are the article's own text line for line. Required context retained uncounted: eq:definition-metric-mercedes (lines 1207--1218). Every definition, display and label of those lines is retained unchanged. Omitted from this package and not redistributed: 1--1206, 1219--1234, 1333--1558. Nothing inside a retained excerpt is omitted or altered: every retained body is delivered exactly as the article's lines are written, so no omission receipt of any kind accompanies this package. Citations: the retained excerpts carry no citation command at all, so no bibliography entry of the article is transcribed and no reference is redistributed. Reference adaptation: none was needed and none was made; every reference inside the delivered unit resolves inside it, so no unresolved reference or citation is produced. Printed slips kept literal: no printed word, symbol, hypothesis or number of the counted statement or of its proof was altered. Layout and class: the article's own class is \texttt{article} with packages including amsmath, amssymb, amsthm, hyperref, algorithm2e, algpseudocode, enumerate, natbib, tikz, pgfplots, adjustbox, subcaption; the wrapper is the lane's pinned \texttt{amsart} editorial wrapper at 10pt on A4, which restates the theorem-like environments the retained text needs and adds the article's own macro definitions that the retained text needs (\texttt{\textbackslash bbR}), with extra wrapper packages (bm, bbm, dsfont, xspace, cancel, mathtools). The delivered numbering is the unit's own. Assets: no retained span contains a figure environment, a graphics-inclusion command, a PDF-page inclusion, a table or a TikZ picture; no image file is copied into this package, the package declares no asset dependency and no image file is redistributed with it. Licence: version arXiv:2607.20677v1 of this article is distributed under the Creative Commons Attribution 4.0 International licence, as recorded in the arXiv licence metadata for this exact version and shown on the abstract page https://arxiv.org/abs/2607.20677v1. A full rights notice is delivered at licenses/arxiv-2607.20677-CC-BY-4.0.txt. Characters: every retained excerpt is pure ASCII, so the delivered engine, XeLaTeX with its default Latin Modern text font, reports no missing character.
\begin{equation}\label{eq:definition-metric-mercedes}
  d((i, x), (j, y)) =
  \left\{\begin{array}{ll}
    |x - y|   & \text{ if } i = j                     \\
    1 - x + y & \text{ if } i = a \text{ and } j = b  \\
    1 - x + y & \text{ if } i = a \text{ and } j = c  \\
    x + 1 - y & \text{ if } i = b \text{ and } j = a  \\
    x + y     & \text{ if } i = b \text{ and } j = c  \\
    x + 1 - y & \text{ if } i = c \text{ and } j = a  \\
    x + y     & \text{ if } i = c \text{ and } j = b. \\
  \end{array}\right.
\end{equation}

\begin{proposition}
  Let $(G, d)$ be the space defined by~\eqref{eq:definition-metric-mercedes}.
  The only functions $f:G \to \mathbb{R}$ that are in $C^1(G)$  are 
  functions with a critical point at $(a, 1) = (b, 0) = (c, 0)$. 
\end{proposition}

\begin{proof}
The claim is proved by showing that the directional derivative in all possible 
directions is zero at the point $(a, 1)$.   

Consider first the geodesics
  $\gamma_{[(a, 0), (a, 1)]}$, $\gamma_{[(b, 1), (b, 0)]}$, $\gamma_{[(c, 1), (c, 0)]}$ and use
  the property of being $C^1$ along geodesics to conclude that there are $C^1$ functions
  $f_\alpha, f_\beta, f_\gamma:[0, 1] \to \bbR$ such that
  \begin{equation*}
    f(i, x) =
    \left\{
      \begin{array}{ll}
        f_\alpha(x) & \text{ if } i = a  \\
        f_\beta(x) & \text{ if } i = b  \\
        f_\gamma(x) & \text{ if } i = c. \\
      \end{array}
    \right.
  \end{equation*}

  The next step is to compute the directional derivatives, expressing
  them using the derivatives of $f_\alpha,f_\beta$ and $f_\gamma$.
This computation yields
  \begin{equation*}
    df(i, x)[\gamma_{[(i, x), (j, y)]}] =
    \left\{
      \begin{array}{ll}
        f_\alpha'(x) (y - x)
        & \!\!\text{if } i = a \text{ and } j = a \text{ and } x \in [0, 1] \\
        f_\alpha'(x) (1 - x + y)
        & \!\!\text{if } i = a \text{ and } j = b \text{ and } x \in [0, 1) \\
        f_\alpha'(x) (1 - x + y)
        & \!\!\text{if } i = a \text{ and } j = c \text{ and } x \in [0, 1) \\
        %
        - f_\beta'(x) (x + 1 - y)
        & \!\!\text{if } i = b \text{ and } j = a \text{ and } x \in (0, 1]  \\
        f_\beta'(x) (y - x)
        & \!\!\text{if } i = b \text{ and } j = b \text{ and } x \in [0, 1] \\
        - f_\beta'(x) (x + y)
        & \!\!\text{if } i = b \text{ and } j = c \text{ and } x \in (0, 1] \\
        %
        - f_\gamma'(x) (x + 1 - y)
        & \!\!\text{if } i = c \text{ and } j = a \text{ and } x \in (0, 1] \\
        - f_\gamma'(x) (x + y)
        & \!\!\text{if } i = c \text{ and } j = b \text{ and } x \in (0, 1] \\
        f_\gamma'(x) (y - x)
        & \!\!\text{if } i = c \text{ and } j = c \text{ and } x \in [0, 1] \\
      \end{array}
    \right. .
  \end{equation*}

  The point $(a, 1) = (b, 0) = (c, 0)$ is of particular interest for the
  topology of its neighborhoods. There are three independent geodesics starting
  from this point, and evaluating the corresponding three directional derivatives yields
  \begin{align*}
    df(a, 1)[\gamma_{[(a, 1), (a, 0)]}] &= f_\alpha'(1) \\
    df(b, 0)[\gamma_{[(b, 0), (b, 1)]}] &= f_\beta'(0) \\
    df(c, 0)[\gamma_{[(c, 0), (c, 1)]}] &= f_\gamma'(0).
  \end{align*}

  In order to obtain the constraints on $f$, we evaluate the directional
  derivatives at points in a neighborhood of the critical point, along geodesics
  passing through this point
  \begin{align*}
    df(a, 1-\varepsilon)[\gamma_{[(a, 1-\varepsilon), (b, 1)]}] &= f_\alpha'(1-\varepsilon) (1 + \varepsilon) \\
    df(a, 1-\varepsilon)[\gamma_{[(a, 1-\varepsilon), (c, 1)]}] &= f_\alpha'(1-\varepsilon) (1 + \varepsilon) \\
    df(b, \varepsilon)[\gamma_{[(b, \varepsilon), (a, 0)]}]     &= - f_\beta'(\varepsilon) (1 + \varepsilon) \\
    df(b, \varepsilon)[\gamma_{[(b, \varepsilon), (c, 1)]}]     &= - f_\beta'(\varepsilon) (1 + \varepsilon) \\
    df(c, \varepsilon)[\gamma_{[(c, \varepsilon), (a, 0)]}]     &= - f_\gamma'(\varepsilon) (1 + \varepsilon) \\
    df(c, \varepsilon)[\gamma_{[(c, \varepsilon), (b, 1)]}]     &= - f_\gamma'(\varepsilon) (1 + \varepsilon).
  \end{align*}

  Since $f$ is $C^1$ along geodesics, we can take the limits as $\varepsilon \downarrow 0$ and
  we obtain the system of equations
  \begin{align*}
    f_\beta'(0) &=   f_\alpha'(1) \\
    f_\gamma'(0) &=   f_\alpha'(1) \\
    f_\alpha'(1) &= - f_\beta'(0) \\
    f_\gamma'(0) &= - f_\beta'(0) \\
    f_\alpha'(1) &= - f_\gamma'(0) \\
    f_\beta'(0) &= - f_\gamma'(0).
  \end{align*}

  This shows that $f_\alpha'(1) = f_\beta'(0) = f_\gamma'(0) = 0$ and that
  \begin{equation*}
    \forall j \in \{a, b, c\}, y \in [0, 1]\quad df(a, 1)[\gamma_{[(a, 1), (j, y)]}] = 0.
  \end{equation*}
  In other words, $f$ has a critical point at $(a,1)$ as claimed. 
%   \begin{equation*}
%     df(a, 1)[{\gamma_{[(a, 1), (a, 0)]}}] =
%     df(a, 1)[{\gamma_{[(a, 1), (b, 1)]}}] =
%     df(a, 1)[{\gamma_{[(a, 1), (c, 1)]}}] = 0.
%   \end{equation*}
\end{proof}

\end{document}
